\section{Discussion and Conclusion}
\label{sec:discussion}
\label{sec:conclusion}

A simple gripper can gain in-grasp dexterity from its workspace rather than from more sensing, actuation, or planning, and the geometry that provides it can be designed in closed form. The method followed here has three steps: identify what passive contacts can and cannot prescribe, choose the fewest constraints that supply each missing motion, and place them by construction. For planar in-grasp manipulation, it leads to two free rods, a segment, and an arc. Because passive contacts cannot determine a motion that combines sliding and turning in the grip, a request is executed as a shift and then a turn. The shift sets the number of rods: two rods hold any admissible direction in one motion, with a closed-form stiffness, where one rod holds only finitely many directions and only indifferently. The turn needs one of them, at an edge or a vertex. The bearing makes the turn's contact wrench independent of rod friction, which yields a friction-free estimate of grasp slip, and across a vertex the second rod separates automatically when the turn begins. A quadratic equation tests whether an installed pair realizes a constructed placement. The individual ingredients are classical; what the paper contributes is a design rule small enough to apply by hand, in which no planning quantity depends on a friction coefficient or a force value. In practice, the tool is two pins on bearings, the gripper and its controller are unchanged, and the manipulation motions use no sensing, which makes the approach a candidate low-cost alternative to a regrasp station or a more dexterous hand. Hardware trials demonstrate turns, including the vertex-contact case; in simulation, shift tilts followed the predicted approach to $e/K_{\mathrm{sh}}$ to within 0.5\%, turn slip followed the first-order estimate to within $0.06^{\circ}$ for $\rho/\ell\leq0.125$, and complete sequences, including one with re-engagement, ended within the predicted slip.

The construction assumes a polygonal outline, convex or not, touched by each rod at a single point, a known interior grasp, a pinch normal to the plane, quasi-static motion, and grip settings that retain the object while allowing sliding and pivoting. The stiffness and slip results are first-order estimates. The slip, and the rate at which a shift approaches the stiffness tilt, rest on an ellipsoidal limit surface for the pad, and the slip also neglects bearing friction; all should be read as design trends rather than exact predictions. The hardware trials used a single rod, retained only mean absolute errors, and cannot test the predicted under-rotation. Fixed spacing, accessible contacts, grasp travel, clearance, and robot reach limit the usable requests; the wrist may rotate before engagement but is held during manipulation. Hardware shifts and complete sequences remain to be tested. Whether one installed pair reaches every usable grasp and orientation by chaining motions, and how to trade the number of motions against the stiffness of each, are still open. So is the extension of the same three-step method to other changes of grasp, such as flipping a part out of the plane, where the missing constraints and their minimal geometry would differ.
